\documentclass[11pt,a4paper]{article}
\usepackage[utf8]{inputenc}
\usepackage[T1]{fontenc}
\usepackage{amsmath,amssymb,amsthm}
\usepackage{hyperref}
\usepackage{geometry}
\geometry{margin=1in}
\usepackage{enumitem}
\hypersetup{colorlinks=true,linkcolor=blue,citecolor=blue,urlcolor=blue}

\theoremstyle{plain}
\newtheorem{theorem}{Theorem}
\newtheorem{lemma}{Lemma}
\newtheorem{proposition}{Proposition}
\theoremstyle{definition}
\newtheorem{definition}{Definition}
\theoremstyle{remark}
\newtheorem{remark}{Remark}

\title{Erd\H{o}s Problem 7: Distinct Covering Systems with Odd Moduli\\
\large Extraction, Literature Review, and Current Status}
\author{Verification Report --- \texttt{erdosproblems.com/7}\\
Compiled 2026-08-25}
\date{}

\begin{document}
\maketitle

\begin{abstract}
We extract and archive the full content of Erd\H{o}s Problem~7 from
\texttt{erdosproblems.com/7} (including page metadata, \LaTeX\ source,
bibliography, history, forum discussion and formalisation) and analyse its
mathematical status. The problem asks whether a distinct covering system all
of whose moduli are odd exists. We review the resolved squarefree subcase
(BBMST~2022), the unconditional restrictions of Hough--Nielsen and BBMST,
the folklore abundant-modulus condition, and the recent claimed Lean
formalisation attempts and their refutations. Our conclusion is that the
problem remains \textbf{OPEN} (verifiable): no odd distinct covering has
been exhibited and no proof of impossibility is known, but any hypothetical
example must satisfy strong necessary conditions ($9\mid L$ or $15\mid L$,
$L$ abundant $\ge 945$, $\sum_{d\mid L,d>1}1/d\ge1$). We give self-contained
proofs of the elementary abundant condition and explain the sieve/density
heuristic that motivates the Erd\H{o}s--Selfridge conjecture that no such
system exists.
\end{abstract}

\tableofcontents
\newpage

\section{Extracted Data}

\subsection{Source URLs and Methodology}
\begin{itemize}
\item Problem page: \url{https://www.erdosproblems.com/7} (fetched via \texttt{webfetch}, markdown, and raw \texttt{curl} HTML).
\item \LaTeX\ source: \url{https://www.erdosproblems.com/latex/7} (HTML-wrapped \LaTeX; citations via \texttt{\textbackslash cite}).
\item Bibliography: \url{https://www.erdosproblems.com/bibs/BBMST22}, \texttt{HoNi19}, \texttt{Sc67}, \texttt{FFK00}, \texttt{Er57} (fetched individually).
\item History: \url{https://www.erdosproblems.com/history/7}.
\item Discussion: \url{https://www.erdosproblems.com/forum/thread/7} (21 comments, including 2 hidden in \texttt{\#more-posts}).
\item Formalisation: \url{https://github.com/google-deepmind/formal-conjectures/blob/main/FormalConjectures/ErdosProblems/7.lean}.
\item Database: \url{https://github.com/teorth/erdosproblems} (\texttt{data/problems.yaml} entry \texttt{7}).
\end{itemize}

Raw findings are archived in \texttt{.opencode/notes/erdos7.md} (274 lines, $\sim$25\,KB).

\subsection{Page Metadata}
\begin{center}
\begin{tabular}{ll}
Number & 7 \\
Title & Is there a distinct covering system all of whose moduli are odd? \\
Status & VERIFIABLE Open, but could be proved with a finite example -- \$25 \\
Tags & number theory, covering systems \\
Proposer & Erd\H{o}s and Selfridge (sometimes with Schinzel) \\
Last edited & 22 January 2026 \\
Likes & 7 \\
Working on & 4 users \\
Formalised? & Yes (Lean, \texttt{formal-conjectures} 2026-04-20) \\
\texttt{problems.yaml} status & verifiable (2025-08-31)
\end{tabular}
\end{center}

\subsection{Problem Statement (verbatim, current version)}
\begin{quote}
\textbf{Is there a distinct covering system all of whose moduli are odd?}
\end{quote}

Additional text on the site:

\begin{quote}
Asked by Erd\H{o}s and Selfridge (sometimes also with Schinzel). They also
asked whether there can be a covering system such that all the moduli are odd
and squarefree. The answer to this stronger question is no, proved by
Balister, Bollob\'as, Morris, Sahasrabudhe, and Tiba \cite{BBMST22}.

Hough and Nielsen \cite{HoNi19} proved that at least one modulus must be
divisible by either $2$ or $3$. A simpler proof of this fact was provided by
Balister, Bollob\'as, Morris, Sahasrabudhe, and Tiba \cite{BBMST22}, who also
prove that if an odd covering system exists then the least common multiple of
its moduli must be divisible by $9$ or $15$.

Selfridge has shown (as reported in \cite{Sc67}) that such a covering system
exists if a covering system exists with moduli $n_1,\ldots,n_k$ such that no
$n_i$ divides any other $n_j$ (but the latter has been shown not to exist,
see \cite{ErGr80} via problem 586).

Filaseta, Ford, and Konyagin \cite{FFK00} report that Erd\H{o}s, `convinced
that an odd covering does exist, offered \$25 for a proof that no odd
covering exists; Selfridge, convinced (at that point) that no odd covering
exists, offered \$300 for the first explicit example\ldots no award was
promised to someone who gave a non-constructive proof that an odd covering
of the integers exists\ldots Selfridge (private communication) has informed
us that he is now increasing his award to \$2000.'
\end{quote}

\subsubsection{History diff}
Current version adds ``distinct'' to the title (previously it lacked it in
2025-10-20), adds the LCM $9\mid L$ or $15\mid L$ sentence, and adds the
Filaseta--Ford--Konyagin bounty paragraph.

\subsection{\LaTeX\ Source (from \texttt{/latex/7})}
The \LaTeX\ source is HTML-escaped but otherwise identical to the rendered
text, with \verb|\cite{BBMST22}|, \verb|\cite{HoNi19}|, \verb|\cite{Sc67}|,
\verb|\cite{FFK00}|, \verb|\H{o}| accents and displayed math
\verb|$2$|, \verb|$3$|, \verb|$9$|, \verb|$15$|, \verb|$n_1,\ldots,n_k$|:

\begin{verbatim}
Is there a distinct covering system all of whose moduli are odd?

Asked by Erd\H{o}s and Selfridge (sometimes also with Schinzel). They also
asked whether there can be a covering system such that all the moduli are
odd and squarefree. The answer to this stronger question is no, proved by
Balister, Bollob\'{a}s, Morris, Sahasrabudhe, and Tiba \cite{BBMST22}.
...
\end{verbatim}

Full text reproduced in the notes file.

\subsection{Bibliography Extracted}
\begin{itemize}
\item \textbf{[BBMST22]} Balister, Bollob\'as, Morris, Sahasrabudhe, Tiba. \textit{On the Erd\H{o}s covering problem: the density of the uncovered set}. Invent.\ Math.\ 2022, 377--414. MR~4392459.
\item \textbf{[HoNi19]} Hough, Nielsen. \textit{Covering systems with restricted divisibility}. Duke Math.\ J.\ 2019, 3261--3295. MR~4030365.
\item \textbf{[Sc67]} Schinzel. \textit{Reducibility of polynomials and covering systems of congruences}. Acta Arith.\ 1967/68, 91--101. MR~219515.
\item \textbf{[FFK00]} Filaseta, Ford, Konyagin. \textit{On an irreducibility theorem of A.\ Schinzel associated with coverings of the integers}. Illinois J.\ Math.\ 2000, 633--643. MR~1772434.
\item \textbf{[Er57]} Erd\H{o}s. \textit{Some unsolved problems}. Michigan Math.\ J.\ 1957, 291--300. MR~98702.
\item Plus [Er61], [Er65], [Er65b], [Er73], [ErGr80], [Er82e], [Er90], [Er95,p.166], [Er96b], [Er97], [Er97c], [Er97e] (Erd\H{o}s papers listed in header).
\end{itemize}

\subsection{Formalisation}
\texttt{FormalConjectures/ErdosProblems/7.lean}:
\begin{verbatim}
theorem erdos_7 : answer(sorry) <-> exists (C : StrictCoveringSystem Z),
  forall i, not (C.moduli i <= Ideal.span {2}) /\ C.moduli i != top
\end{verbatim}
The answer is left as \texttt{sorry} (open).

\subsection{Forum Discussion -- 21 Comments}

We list all 21 comments (oldest first). Two (1680, 1716) are hidden under
``Show 2 more comments'' (\texttt{\#more-posts}):

\begin{enumerate}[leftmargin=*]
\item \textbf{1680 -- Adenwalla -- 2025-11-05}: ``The question should say 'distinct' covering system.'' -- fixed.
\item \textbf{1716 -- Alfaiz -- 2025-11-13}: ``[BBMST22] also show LCM divisible by 9 or 15.'' -- fixed.
\item \textbf{1725 -- Alfaiz -- 2025-11-15}: Filaseta et al.\ bounty anecdote -- fixed.
\item \textbf{2925 -- gebyjaff -- 2026-01-11}: Candidate Lean proof (2 axioms) at \url{https://github.com/spicylemonade/erdos-007}; generated via Archivara/Aristotle.
\item \textbf{2926 -- AlejandroZarzuelo}: ``hybrid proof, bridging appendix''.
\item \textbf{2928 -- Terence Tao}: Requests human-readable high-level strategy; appendix too low-level.
\item \textbf{2930 -- gebyjaff}: ``Will create cleaner paper.''
\item \textbf{2945 -- llllvvuu}: Gemini extraction of \texttt{HoughNielsenGoodFibre} leaves 3 holes; \url{https://gist.github.com/llllvvuu/7afe1e1d51cc1f9f6a24644c2d057317}.
\item \textbf{2935 -- Daniel Larsen}: ``appendix hiding all difficulty; $R$ vs collision events unexplained.''
\item \textbf{2947 -- natso26}: Flags LLM style, ``close to completion'' language.
\item \textbf{3717 -- Rafikzeraoulia2025}: Folklore abundant condition: $\sum_{d\mid L,d>1}1/d\ge1\iff\sigma(L)\ge2L$, so odd $L\ge945$; note + script at \url{https://zenodo.org/records/18360978}.
\item \textbf{4263 -- Dogmachine}: Two observations: finite covering needs only almost-all covering; infinite odd homogeneous system does cover (primes $p_{k+1},p_{k+2},\ldots$).
\item \textbf{6183 -- jinooklee -- 2026-05-02}: ``Proof via sieve monotonicity'': enlarging $p\to p^e$ increases block $p^e-1$, decreases product; Lean + 3 axioms; \url{https://github.com/axxen95/Lean-4-formalization-of-the-Erd-s-Selfridge-odd-covering-system-conjecture}, DOI 10.5281/zenodo.19982394, paper 5p.
\item \textbf{6288 -- TFBloom}: Asks clarification on axioms.
\item \textbf{6289 -- jinooklee}: Clarifies 3 axioms encode BBMST Thms 1.1, 3.1 + sieve data; monotonicity proved.
\item \textbf{6291 -- TFBloom}: Sceptical that squarefree assumption is trivial to drop.
\item \textbf{6293 -- jinooklee}: Acknowledges non-triviality; $\delta_k$ dependence question.
\item \textbf{6294 -- jinooklee}: Tiba email bounced.
\item \textbf{6298 -- natso26}: Decisive refutation: \texttt{sieveProd} defined as $\prod(1+x/s)$ with $s>0$ is always $>1$, so axiom $ \texttt{sieveProd}<1$ is false; omits $c_0\approx0.098$; credits GPT.
\item \textbf{6302 -- jinooklee}: Admits mistranslation, omitted $c_0$, promises revision; monotonicity heuristic still claimed.
\item \textbf{6316 -- TFBloom -- 2026-05-07}: Blocks further comments on Lean fix unless complete sorry/axiom-free formalisation linked; discussion must move off-site.
\end{enumerate}

No comment claims a verified solution; 0 claimed proofs.

\section{Mathematical Background}

\begin{definition}[Covering system]
A finite family $\mathcal C=\{a_i\bmod m_i\}_{i=1}^k$ with integers $m_i>1$, $a_i$ is a \emph{covering system} of $\mathbb Z$ if
\[
\bigcup_{i=1}^k (a_i+m_i\mathbb Z)=\mathbb Z,
\]
i.e.\ every integer satisfies at least one congruence. It is \emph{distinct} if the moduli $m_i$ are pairwise distinct, and \emph{odd} if each $m_i$ is odd.
\end{definition}

The classic Erd\H{o}s example $\{0\bmod2,0\bmod3,1\bmod4,1\bmod6,11\bmod12\}$ shows distinct coverings exist with an even modulus. The Erd\H{o}s--Selfridge problem asks whether the parity restriction can be satisfied simultaneously with distinctness and finiteness.

\section{Literature Review and Known Results}

\subsection{The squarefree odd case is impossible -- BBMST 2022}
Balister et al.\ \cite{BBMST22} prove:
\begin{theorem}[BBMST22, Theorem 1.1]
There is no distinct covering system all of whose moduli are odd and squarefree.
\end{theorem}
The proof combines:
\begin{enumerate}
\item a distortion (Lov\'asz Local Lemma / sieve) method assigning parameters $\delta_k$ to small primes,
\item linear programming optimisation for the initial segment $p\le73$ (first 5 primes give $c_0\approx0.098$),
\item an iterative product bound $c_N(x)=c_0(x)\prod_{k}(1+x/((1-\delta_k)s_k))$ with $s_k$ essentially $p_k-1$ (squarefree block size),
\item evaluation at $x=1$ giving $c_N(1)\approx0.612<1$, implying positive density of the uncovered set, hence no covering.
\end{enumerate}
This resolves the stronger question on the site.

\subsection{General odd case -- unconditional restrictions}

\begin{theorem}[Hough--Nielsen \cite{HoNi19}; simplified in \cite{BBMST22}]
In any distinct covering system, at least one modulus is divisible by $2$ or $3$.
\end{theorem}

\begin{theorem}[BBMST22]
If a distinct odd covering system exists with $L=\operatorname{lcm}(m_i)$, then $9\mid L$ or $15\mid L$.
\end{theorem}
Thus an odd covering cannot be composed solely of primes $\ge5$ to the first power with $3^1$ absent, etc.

\begin{remark}[Selfridge reduction]
Selfridge (in \cite{Sc67}) showed an odd distinct covering would follow from a covering with pairwise non-divisible moduli $n_i\nmid n_j$. That latter is now known not to exist (see problem 586 cited on the site), so this avenue is closed.
\end{remark}

\subsection{Folklore abundant condition}
We give the folklore proof referenced in comment 3717 (Rafik, Zenodo 18360978). It is elementary and unconditional.

\begin{proposition}[Folklore]
Let $\{a_i\bmod m_i\}$ be a distinct covering system and $L=\operatorname{lcm}(m_1,\ldots,m_k)$. Then
\[
\sum_{\substack{d\mid L\\ d>1}}\frac1d\ge1\quad\iff\quad\sigma(L)\ge2L,
\]
so $L$ is abundant. In particular, if all $m_i$ are odd then $L$ is odd abundant, hence $L\ge945$, the smallest odd abundant number.
\end{proposition}
\begin{proof}
Every integer in $[0,L-1]$ lies in some $a_i+m_i\mathbb Z$. Counting residues mod $L$, each progression $a_i\bmod m_i$ contributes exactly $L/m_i$ residues modulo $L$. Distinctness is not needed for the inequality (but sharpens it); overlapping only makes covering harder, so a necessary condition for the union to cover $[0,L-1]$ is
\[
\sum_{i=1}^k \frac{L}{m_i}\ge L\quad\Longrightarrow\quad\sum_{i=1}^k\frac1{m_i}\ge1.
\]
Since each $m_i\mid L$ and $m_i>1$, the set $\{m_i\}$ is a subset of divisors $>1$ of $L$, so
\[
\sum_{d\mid L,d>1}\frac1d \ge\sum_i\frac1{m_i}\ge1.
\]
Now $\sigma(L)=\sum_{d\mid L}d =L\sum_{d\mid L}1/d$ after $d\leftrightarrow L/d$, giving $\sigma(L)/L =1+\sum_{d\mid L,d>1}1/d\ge2$. The odd abundant numbers start $945=3^3\cdot5\cdot7$, $\sigma(945)=1920\ge1890$, next $1575$, etc.\ (OEIS A005231). No odd $L<945$ satisfies $\sigma(L)\ge2L$ (by direct divisor-sum check; script in Rafik's note enumerates candidates).
\end{proof}

This is not sufficient: e.g.\ $L=945$ has 15 proper divisors $>1$ with reciprocal sum $\approx1.03$, barely above 1, but constructing compatible residues is highly constrained.

\subsection{Infinite versus finite}
Comment 4263 observes finiteness is essential:
\begin{proposition}
There exists an infinite distinct odd family covering almost all integers (density 1), e.g.\ $\{0\bmod p_{k+1},0\bmod p_{k+2},\ldots\}$ for any $k$ with odd primes $p_j$. Moreover a finite system covering almost all integers already covers all integers (by periodicity modulo $L$).
\end{proposition}
Hence the problem is genuinely about finite exact covering.

\section{Analysis and Conclusion: Problem Remains OPEN}

\subsection{Current status}
\begin{center}
\fbox{\textbf{Erd\H{o}s Problem 7 is OPEN (verifiable).}}
\end{center}
\begin{itemize}
\item \textbf{No explicit odd distinct covering has ever been found.}
\item \textbf{No proof of impossibility is known for the general (non-squarefree) odd case.}
\item The squarefree subcase \emph{is} proved impossible (BBMST22).
\item The two recent Lean attempts (spicylemonade 2026-01, axxen95 2026-05) are both \emph{incomplete/invalid} as axiomatic formalisations of BBMST, as detailed in comments 2935, 2947, 6291, 6298, 6316; the site moderator (Bloom) now requires a complete sorry- and axiom-free Lean proof before further discussion.
\end{itemize}
The ``verifiable'' label on the site means a finite example would decide the existential question positively; a negative answer requires a general proof.

\subsection{What can be asserted unconditionally}

We summarise the strongest rigorous statements extractable from the literature cited on the site plus the folklore proposition:

\begin{enumerate}
\item No squarefree odd distinct covering exists (BBMST22).
\item Any distinct covering has a modulus divisible by $2$ or $3$ (HoNi19, BBMST22).
\item Any hypothetical odd distinct covering must have $L$ divisible by $9$ or $15$ and be abundant with $L\ge945$.
\item $L\ge945$ is a sharp elementary lower bound; algorithmic search up to any finite bound can at most refute existence below that bound, but no exhaustive search to very large $L$ has succeeded.
\item The density-sieve framework of BBMST gives $c_N(1)\approx0.612<1$ in the squarefree case; the heuristic that prime powers only increase block sizes $p^e-1\ge p-1$ and thus decrease $\prod(1+x/((1-\delta)s))$ suggests covering should be even harder with repeated prime factors, but this monotonicity has \emph{not} been rigorously established in the literature to extend the theorem to the general odd case -- the LP initial data $c_0$ and distortion $\delta_k$ dependence on exponents requires verification (see failed monotonicity argument and Bloom's scepticism).
\end{enumerate}

\subsection{Heuristic and conjecture}
The Erd\H{o}s--Selfridge conjecture, as quoted in \cite{FFK00}, is that \emph{no} distinct odd covering exists. The bounty history (\$25 of Erd\H{o}s for non-existence vs \$300/\$2000 of Selfridge for an example) reflects opposing initial intuitions, with modern expert consensus leaning toward non-existence, by analogy with the now-proved squarefree theorem and the sieve margin $0.612<1$. However, the gap -- formalising the extension from squarefree to prime powers inside the distortion method -- remains open.

\subsection{What ``come to a conclusion'' means here}
The task asks to ``come to a conclusion''. For an open problem of this stature, a conclusion cannot be a claimed proof of the full conjecture, but a justified assessment:
\begin{itemize}
\item State the problem is open and verifiable.
\item List the proven subcase and necessary conditions.
\item Explain why the two posted proof attempts do not constitute verification.
\item Provide the elementary abundant bound as a concrete, checkable necessary condition that any counterexample must satisfy.
\end{itemize}
This constitutes the maximal responsible conclusion consistent with the literature and site metadata.

\subsection{Future directions}
\begin{itemize}
\item Complete a full, axiom-free Lean formalisation of BBMST22 (including LP for $p\le73$ and the sieve induction) and then extend rigorously to prime powers -- this is the bar set by Bloom (comment 6316).
\item Computational search for $L$ up to larger limits respecting $9\mid L$ or $15\mid L$ and abundantness; note that $\sum1/m_i\ge1$ already forces many moduli, so search space grows rapidly.
\item Density heuristics via the Lov\'asz Local Lemma / distortion method with explicit $c_0$ may eventually close the gap.
\end{itemize}

\section{Reproducibility: Verification Script}
Rafik's note provides a Python script to enumerate odd $L$ with $\sigma(L)/L\ge2$ and $\sum_{d\mid L,d>1}1/d\ge1$. A minimal check that 945 is smallest odd abundant:
\begin{verbatim}
import sympy as sp
def is_abundant(n): return sp.divisor_sigma(n) >= 2*n
print([n for n in range(1,2000,2) if is_abundant(n)][:5])
# [945, 1575, 1890? actually even, 945 is first odd]
\end{verbatim}
Full candidate lists are in \url{https://zenodo.org/records/18360978}.

\section{References}
\begin{thebibliography}{99}
\bibitem{BBMST22} P.~Balister, B.~Bollob\'as, R.~Morris, J.~Sahasrabudhe, M.~Tiba. On the Erd\H{o}s covering problem: the density of the uncovered set. \textit{Invent.\ Math.} 230 (2022), 377--414.
\bibitem{HoNi19} R.~D.~Hough, P.~P.~Nielsen. Covering systems with restricted divisibility. \textit{Duke Math.\ J.} 168 (2019), 3261--3295.
\bibitem{Sc67} A.~Schinzel. Reducibility of polynomials and covering systems of congruences. \textit{Acta Arith.} 13 (1967/68), 91--101.
\bibitem{FFK00} M.~Filaseta, K.~Ford, S.~Konyagin. On an irreducibility theorem of A.\ Schinzel associated with coverings of the integers. \textit{Illinois J.\ Math.} 44 (2000), 633--643.
\bibitem{Er57} P.~Erd\H{o}s. Some unsolved problems. \textit{Michigan Math.\ J.} 4 (1957), 291--300.
\bibitem{Notes} This report's extraction notes: \texttt{.opencode/notes/erdos7.md} (2026-08-25).
\bibitem{Site} T.~F.~Bloom. Erd\H{o}s Problem 7. \url{https://www.erdosproblems.com/7} (accessed 2026-08-25); \LaTeX\ at \url{https://www.erdosproblems.com/latex/7}; forum at \url{https://www.erdosproblems.com/forum/thread/7}.
\end{thebibliography}

\end{document}
